% Tree / hierarchy diagram template (forest)
% Compile with: pdflatex tree-template.tex
\documentclass[border=10pt]{standalone}
\usepackage[T1]{fontenc}
\usepackage{forest}
\usepackage{xcolor}

% colors
\definecolor{leaffill}{RGB}{220, 233, 245}
\definecolor{leafborder}{RGB}{120, 165, 210}
\definecolor{boxborder}{RGB}{70, 70, 70}

% per-tier styles (give each level a fixed text width to align columns)
\tikzset{
  rootbox/.style={
    draw=boxborder, rounded corners=3pt,
    text width=1.7cm, align=center,
    font=\footnotesize, inner sep=4pt, minimum height=10mm,
  },
  l1box/.style={
    draw=boxborder, rounded corners=3pt,
    text width=3.4cm, align=center,
    font=\footnotesize, inner sep=5pt, minimum height=7mm,
  },
  catbox/.style={
    draw=boxborder, rounded corners=3pt,
    text width=3.4cm, align=center,
    font=\footnotesize, inner sep=5pt, minimum height=7mm,
  },
  leafbox/.style={
    draw=leafborder, line width=0.5pt, fill=leaffill,
    rounded corners=3pt,
    text width=4.8cm, align=left,
    font=\scriptsize, inner sep=4pt,
  },
}

\begin{document}
\begin{forest}
  for tree={
    grow'=east,
    parent anchor=east, child anchor=west,
    l sep=14pt, s sep=7pt,
    edge={-, thin, draw=boxborder},
    % L-shaped connector with rounded bend
    edge path={
      \noexpand\path[\forestoption{edge}, rounded corners=4pt]
        (!u.parent anchor) -- +(8pt,0) |- (.child anchor);
    },
  },
  where level=0{rootbox}{},
  where level=1{l1box}{},
  where level=2{catbox}{},
  where n children=0{leafbox}{},
  [Root
    [Category A
      [Sub A.1 [{Leaf content with as many wrapped lines as needed.}]]
      [Sub A.2 [{Another leaf.}]]
    ]
    [Category B
      [Sub B.1 [{Leaf content.}]]
    ]
  ]
\end{forest}
\end{document}
